\begin{established}[The real row and measured columns at the prescribed scale]\label{hyp:real-row-measured-models}
Fix an arithmetic row specialization of Proposition~\ref{prop:joint-row-orders}. Let \(F_u^0\) be
its degree-zero column form model for the word \(u\), including
the chosen column values, the level at \(T_u\) and the two
type markings. If \(V_{\rm good}\subset V_{\rm full}\) are the two rows recorded here, the dimension conditions used in the replacements are
\[
 \dim F_u^0>0\quad\text{for each word }u,\qquad\dim V_{\rm good}\le\dim V_{\rm full}.
\]
The local matrix models, Petersson measure and quotient measures on the flag fibres can all be computed in the same Hermitian coordinates. In these coordinates the unit lowest vectors satisfy the absolute Bessel bounds used for the finite term expansion.
\end{established}
\begin{proof}
The dimension formula of Proposition~\ref{prop:marked-column-dimension} is a product of positive volume, weight and finite marking factors. The other dimension inequality follows from inclusion of vector spaces. Keep this row and these columns while making the following two
replacements of the spectral data. In the first replacement use
the scale of the actual Haar data of the form, the discrete-series
parameters obtained from the labelled Harish--Chandra weights
and signatures, and the numerical discriminant of the fixed
scale context.

For each finite place the local matrix model is
obtained from the completion of the marked quadratic algebra
with its involution. The adjoint form in the local matrix model comes from scalar extension of the same Hermitian form. For a Gram matrix \(H\) in the row or column coordinates, the adjoint is read locally by
\[
 A\longmapsto H^{-1}\bar A^{\,t}H ,
\] This gives the placewise
isomorphisms for the replaced spectral data, with the same
exponent and size context. Division lattices, their height
functions, the chosen character, and the localization grid
are transported by the row identifications. The transport for
the initial state carries the origin comparison, and that for
the rebuilt state carries the origin and geometric norm
comparisons.
\label{cor:real-row-realization-1}

Let \(X\) be the adelic unitary model of this Hermitian space.
The measures in \(X\) are chosen together with the standard
additive character. Start with the self-dual lattice from
the reference form \(F_{u_*}^0\). At the marked split primes keep the fixed frames. For each finite place \(v\), take the self-dual Haar measure on the additive unipotent coordinates for the standard character \(\psi_v\). In the Fourier convention this means
\[
 \widehat{\widehat f}(x)=f(-x)\qquad(f\text{ a Schwartz--Bruhat function}).
\]

At the real places use the prescribed unit
lowest vectors and the Bessel functions of the discrete
series of the selected weight. The local choice is made
for the weight gaps of the scale context. The stabilizer
quotient measures come from the Gram fibres. Use the
convergent unitary Siegel--Weil identity for the two
stabilizers and the rational quotient density to normalize
their covolumes simultaneously. In this normalization the
Petersson measure and the product measures on the flag
fibres refer to the same ambient Haar, as in the fibre
comparison of Section~\ref{sec:absolute}. The gap bound
gives the weighted absolute integrability of the Bessel
functions through the matrix-coefficient estimate of Proposition~\ref{prop:bessel-sign-gap-decay} and the local Bessel projection in Section~\ref{sec:absolute}. These are the measure and decay
conditions used in the finite term expansion.
\label{cor:real-row-realization-2}
\end{proof}

\begin{proposition}[Identification of the measured columns and their fibres]\label{prop:measured-columns}
For the Hermitian column models and measure of Lemma~\ref{hyp:real-row-measured-models}, match the ordered
frames, the distinguished line of norm \(-T_u\), and the characters
at both type primes. The resulting finite adelic transport induces an
isometry of the degree-zero column spaces. For a nonempty term fibre \(\mathcal F\), write \(\mathcal F/U\) for its lift classes modulo flag unipotents, \(S\) for the marked target stabilizer and \(X_0\) for the target isotropic space. There is an injection
\[
 \mathcal F/U\lhook\joinrel\longrightarrow\operatorname{im}\bigl(S\to\operatorname{Aut}(X_0^\perp/X_0)\bigr),
\]
with finite target. Corank and the existence of a scalar lift are
independent of the representatives of a pair with equal row reading.
\end{proposition}

\begin{proof}
\proofstep{Matching the local lattices.}
For each sequence, two ordered frames have the same Gram matrix.
Their hyperbolic pairs have pairing one, and the distinguished
line has norm \(-T_u\). If \(A\) is the matching between their spans and \(G,G'\)
are their Gram matrices in any bases, the Gram equalities give
\(A^*G'A=G\). This extends by Witt's theorem. In the
orthogonal complement the two degree-zero lattices are
self-dual off \(T_u\). At an unramified inert place split off a unit-norm line at each step. If a chosen vector has unit norm \(a\), surjectivity of the norm on units gives a unit \(c\) with
\[
 N(c)=a^{-1},\qquad \langle cv,cv\rangle=N(c)a=1 .
\] At a split place the two halves of a self-dual
lattice are paired duals. At \(T_u\) the dual quotient on
each split factor has length one. Indeed the primitive
distinguished line splits rationally from its orthogonal
complement in a self-dual middle lattice.

Projection gives
the two dual quotients the same length, and the norm
\(-T_u\) gives that length at the split factors. In split
coordinates the resulting lattice chain is a full lattice
and the inverse image of a hyperplane modulo a uniformizer.
A change of basis \(A\) on one factor forces the inverse transpose on the dual factor, since for paired coordinate vectors
\[
 (Ax)^t(A^{-t}y)=x^ty .
\] At the two type
primes a diagonal unit correction in the common finite
reduction matches the superdiagonal characters in the
chosen frames. Conjugation by the inverse of
\(\operatorname{diag}(1,b,b^2,\ldots)\) multiplies all
first-superdiagonal entries by \(b\). In reversed dual
coordinates inverse transpose negates their sum. The
primitive residue characters are related by a nonzero
residue scalar, so one chooses \(b\) to match the inverse
character in the sharp reading and lifts its entries to units.

\proofstep{Transport of the class sets.}
These local comparisons are identity at almost all places.
They give a finite adelic transport identifying the
column levels and compact models. For the induced bijection
\(a\mapsto a'\) on double quotients, let \(\mathcal I_a\) be the
unitary value comparison in the compact model at \(a\). It sends a
class function \(F\) to \(F'\) with \(F'(a')=\mathcal I_a F(a)\).
Covariance extends this identity to all representatives, and
\[
          \sum_{a'}\|F'(a')\|^2=\sum_a\|F(a)\|^2
\]
uses the bijection and the finer-level trivial stabilizers. The chosen unit vectors in the compact
highest-weight representations identify the invariant
norms. At the finer level the stabilizers are trivial by
the split type-prime argument of
Section~\ref{subsec:selected-word-level}. Indeed the class bijection preserves the values in the identified compact representations. With trivial stabilizers the norm comparison is simply reindexing: for a column function \(f\) and class bijection \(b\),
\[
 \sum_C\|f(C)\|^2=\sum_{C'}\|f(b^{-1}C')\|^2 .
\] It identifies the actual column forms
entering the Bessel and Gram computations.

The finite
term fibres use these very classes and their shorter
and complete flag representatives. Fix one lift \(s_0\) of a term in a nonempty fibre. Write
\(X_0\) for the isotropic space of its marked target flag and
\(\Lambda'\) for its target lattice. Modulo the two flag unipotents the transition for any other lift is \(s_0 c\). If \(b_i\) are the marked flag basis vectors and \(v\) the marked vector, then \(c\) satisfies
\[
 c\Lambda'=\Lambda',\qquad c b_i=b_i\ (\text{each }i),\qquad c v=v .
\] Choose one such \(c\) per lift. Its
action on \(X_0^\perp/X_0\) lies in the finite image of the marked
lattice stabilizer. Here the definite lattice automorphism bound
applies with the signature at stage zero and the positive norm
parameter \(T_{u'}\) of the target word. Write \(\bar c\) for the induced endomorphism of
\(X_0^\perp/X_0\). If \(\bar c=\bar c'\), put \(u=c^{-1}c'\).

The three conditions defining the flag-unipotent equivalence are
\[
 u|_{X_0}=1,\qquad u\Lambda'=\Lambda',\qquad
 (u-1)X_0^\perp\subseteq X_0 .
\]
The last inclusion is just \(\bar c^{-1}\bar c'=1\); the first
two follow from the marking and the integral lattice. Such \(c\) also normalizes
the target flag unipotents, so the two lifts represent the same
class. The chosen map into the finite image is injective.

\proofstep{Independence of representatives.}
Consider also a pair of lifts having the same row reading, and a
second pair with the same respective terms and again equal row
readings. In the transport just used, equality of the two marked
embedding readings identifies the induced Levi actions. The neat
level makes their rational stabilizer representatives equal. In particular the relative target transporter induces the identity on the flag quotient:
\[
 (c_2c_1^{-1})v-v\in X_0\qquad(v\in X_0^\perp).
\] Conjugate it by the first lift. The new relative
matrix differs from the old one by this conjugate and the
lattice-preserving source and frame unipotents. They carry the
space defining the scalar \(s_0(\cdot,\cdot)\) isomorphically
onto the new space, preserving its codimension, and carry any
scalar lift to a scalar lift. Interchange the two pairs to get
the reverse implication. This proves representative independence
before using absolute convergence to reindex the expansion.
\cite{companion-row-count}\label{cor:real-row-realization-3}
\end{proof}

The product discrete series exists with the prescribed
labelled real parameters. Its formal degrees are computed
in the actual product Haar. In the real tie their product
is the one at the true scale, and the finite factor is the
value of the actual counting projector at the identity.
The comparison of quotient volume with mass unfolds the
same product measures and uses the definite genus
decomposition. Hence the identity contribution in the
measured model is exactly the predicted \(D_*\) of
Section~\ref{sec:count}. In particular it uses the
finite group quotient and level index of this row. The
finite good operator is fixed under the conjugate reading
in the matched frames. The local series projector has
its coefficients paired by dualizing the supported tame
labels and reversing their order. This is the coefficient
identity which compares the two readings of the operator.
\label{cor:real-row-realization-4}

\begin{proposition}[The isometry of the refined row]\label{prop:refined-real-row}
Let \(J\) be the sharp geometric level, \(n_P\) its level index
in the normalization of the Petersson form, and \(M(J)\) the full
multiplicity space. Evaluation at the chosen unit lowest vectors
identifies \(M(J)\) with the conjugate measured row, with one positive
squared-norm ratio \(c\). It respects the good projector. If \(j\) is the normalized composite on the arithmetic good row and \(V_{\rm full}\) the row before projection, then
\[
 \|jv\|^2=\|v\|^2\quad(v\text{ in the arithmetic good row}),\qquad \dim V_{\rm full}=D_{\mathrm{all}}.
\]
\end{proposition}

\begin{proof}
\proofstep{The two finite characters.}
Let \(J\) be the geometric level defining the refined row space and \(K_2\) the
corresponding subgroup in \(X\). Use the simultaneous
finite reduction of the type frames to first match
principal kernels. In the common finite group the two
marked sharp levels are conjugate. Lift the finite
conjugator inside the compact group and correct the
finite translate. If \(\iota\) is the group identification,
the corrected translate \(t\) satisfies
\[
 t^{-1}\iota^{-1}(K_2)t=J,\qquad
 \chi_J(t^{-1}\iota^{-1}(k)t)=\overline{\chi_X(k)}
 \quad(k\in K_2).
\]
\proofstep{Inverting lowest-vector evaluation.}
The central idempotents are unchanged by this conjugation
of frames. The real representation comparison uses the
contragredient discrete series at the same infinitesimal
character in the conjugate Hilbert space. Choose the
geometric unit lowest vectors corresponding to those
of \(X\). For an intertwiner in the full sharp multiplicity
space \(M(J)\), evaluation at these lowest vectors after
the translate \(t\) defines a map \(L\) onto the conjugate
row space of \(X\). Automorphy and the displayed level
and character conditions give the prescribed finite
equivariance. At infinity, for a unit lowest vector \(w\)
in a discrete series, let \(P_{v,w}\) be the operator
of its normalized matrix coefficient taking the column
\(w\) to \(v\). Schur orthogonality gives the rank-one matrix-unit identities
\[
 P_{w,w}^2=P_{w,w},\qquad P_{v,w}P_{w,w}=P_{v,w}.
\]
A row function \(f\) in the image of \(P_{w,w}\) gives the
intertwiner \(v\mapsto P_{v,w}f\) by covariance under
real translation of the matrix coefficients.

Evaluation
at \(w\) recovers \(f\). Apply the product of these
operators, using the contragredient identification and
the inverse finite translation. This gives surjectivity
and the inverse to evaluation. The construction goes
through the smooth discrete-series vectors, using the
matrix-coefficient projection with the actual formal
degrees. Rapid decay of cusp forms makes the evaluations
square integrable.
Petersson integration on the two identified quotients has
a single positive ratio \(c\). With the real vectors of
norm one this gives
\[
 \|LS\|^2=c\|S\|^2\qquad(S\in M(J)).
\]
The map is complex linear into the space of conjugated
functions with its complex structure there. The
finite action on both sides now intertwines with
conjugation of its coefficients. Thus \(L\) carries the
good part of \(M(J)\) onto the conjugate good subspace,
as seen by applying the central good projector inside
the full evaluation isomorphism.
\label{cor:real-row-realization-5}

\proofstep{Normalization on the good subspace.}
The geometric realization of the arithmetic row
\(E_{\rm row}\) is an isomorphism onto the linearized
good multiplicity space at \(J\). In the multiplicity evaluation, write \(d\mu\) for the measure before sharp-level normalization. The Petersson measure on the row is
\[
 d\mu_{\rm sharp}=n_P^{-1}d\mu .
\]
The good space at \(J\) is contained in the principal
model. Invert the faithful principal realization on
this subspace to obtain the geometric row embedding.
Its composition \(e_E\) with multiplicity realization
has image the whole good space of \(M(J)\), by surjectivity
of this geometric realization. Using the same unit tuple
in both evaluations gives \(\|e_E x\|^2=n_P\|x\|^2\).

It follows that
\[
             j=(c n_P)^{-1/2}L e_E
\]
is the isometry onto the conjugate good row space of the
measured model. The full evaluation bijection identifies
its full dimension with the recorded \(D_{\mathrm{all}}\).
This dimension is from the full space at the refined level before applying
the good projector. The two positive factors in \(j\) cancel
exactly in the pairing for every host and every allowed
word level.
\label{cor:real-row-realization-6}
\end{proof}

For the second replacement choose the spectral datum
using the row with its normalized Petersson measure and the comparison map \(j\). This choice fixes
\(t,c,j\) before any replacement of the division lattices
and heights on the word spaces. If a replacement retains
the degree-zero levels and readings, the term fibres are
carried by the two induced bijections of class labels.
The fibre of a term is the image of its original finite
fibre under those bijections, with equal function values.
Write \(b\) for the isometry induced by reindexing the good space. The transported row map is the composite, so for every row vector \(v\),
\[
 j_{\rm new}=b\circ j,\qquad\|j_{\rm new}v\|=\|jv\|=\|v\|.
\]

To apply
this comparison here take the original division lattices
and their height function. Transport first through the
symmetrized column form model, which preserves the
degree-zero values, then through the equivalence to the
actual rebuilt carriers of the same columns. It follows
that the new comparison of the row spaces at the refined level uses the original
column levels and actual column forms. The full
dictionary, the real tie, the discriminant and the
transported order and character data constitute the
level-four output for this pass.
\label{cor:real-row-realization-7}

\begin{lemma}[The ambient lattices in the used tuples]\label{lem:real-ambient-tuples}
At the sharp and bottom levels separately, inclusion of the used
tuples in the full tuple list identifies their coefficient spaces,
integral ambient lattices and positive Hermitian forms. Writing \(b\) for this identification at one of the two levels, \(\Lambda\) for the lattice, and \(H\) for the form, one has at each used tuple \(\mu\)
\[
 b\Lambda_{\mu,{\rm used}}=\Lambda_{\mu,{\rm full}},\qquad b^*H_{\mu,{\rm full}}=H_{\mu,{\rm used}}
\]
at every embedding where the form is positive. These identifications use the
same degree-zero columns as Proposition~\ref{prop:refined-real-row}.
\end{lemma}

\begin{proof}
The ambient lattices at the sharp and bottom levels use
the rational coordinates of this row specialization. Fix one of these levels and a tuple \(\mu\). Let
\(V_\mu\) be its coefficient-field coordinate space, with
lattice \(\Lambda_\mu\) and Hermitian matrix \(H_\mu\).
Scalar extension at the marked embedding sends its coordinate vectors to a
basis at the marked embedding, and carries \(\Lambda_\mu\)
onto the marked ambient lattice. At a positive embedding
\(\sigma\), choose an orthonormal frame for
\(\sigma(H_\mu)\).

Let \(F\) send the marked basis coordinates to coordinates in this orthonormal frame at \(\sigma\). Thus \(F^*F=\sigma(H_\mu)\), and
\[
 \|F\sigma(x)\|^2=\sigma(x)^*F^*F\sigma(x)=\sigma(x)^*\sigma(H_\mu)\sigma(x).
\] This is the positive form
in the comparison of the ambient norms. At the marked
embedding the original Hermitian and operator identities
give the designated complex form and adjoint under the
reading. Finite generation and spanning of the ambient
lattice follow from those of \(\Lambda_\mu\) and the two
identities for the basis and lattice.

The used tuple list maps into the full tuple list by the
same tuple of local character labels. The prescribed tuple
map agrees with this inclusion because the row indices are
defined by those labels. At either level the coordinate map on the selected tuples and the restriction of the full coordinate map agree on every basis vector. If these maps are \(f,g\),
\[
 f\!\left(\sum_i a_i e_i\right)=\sum_i a_i f(e_i)=\sum_i a_i g(e_i)=g\!\left(\sum_i a_i e_i\right).
\] The identity on coordinates, with the
dimension identified through that equality of indices,
therefore identifies the rational readings, integral
lattices and positive forms at every embedding where the
form is positive. This constructs the two ambient
identifications consumed by the division construction.
The sharp and bottom matrices and lattices are taken from
their respective levels.
\label{cor:real-row-realization-8}
\end{proof}

Both passes have this construction at their own
level-three states. The first level-four state is
completed to a built member before taking its signed
export. The rebuilt pass makes the corresponding choices
on the carried specialization. Their division and theta
constructions consume the corresponding real isometry and
ambient identifications on that pass.
\label{cor:real-row-realization-9}
